\documentclass[10pt]{article}
\usepackage[margin=20mm]{geometry}
\usepackage{amsmath,amssymb,mathtools,fontspec,unicode-math,xurl,hyperref}
\IfFontExistsTF{DejaVu Serif}{\setmainfont{DejaVu Serif}}{\setmainfont{Latin Modern Roman}}
\setmathfont{latinmodern-math.otf}
\setlength{\parindent}{0pt}\setlength{\parskip}{6pt}\sloppy
\title{Technical review of
The non-Hermitian skin effect with three-dimensional long-range coupling}
\author{AI-assisted technical review}
\date{}
\begin{document}
\maketitle
Abstract: This initial review concerns the 39-page EMS online-first article of 4 August 2025 by Ammari, Barandun, Cao, Davies, Hiltunen and Liu (DOI: 10.4171/JEMS/1685), supplied through CaltechAUTHORS. One local error is confirmed in the spectral inclusion in Theorem A.3. A unilateral-shift example disproves the inclusion at the stated exponential scale. The main localization argument instead invokes the separately proved Theorem A.4; no damage to that argument is established by this example. The assessment is Minor within the coverage specified below.\par
\section*{Scope and principal dependencies}
The paper relates subwavelength resonances to a gauge capacitance matrix, then establishes localized pseudo-eigenvectors using a banded approximation. The checks covered the physical model and conventions, Theorem 4.4, the Toeplitz approximation and off-diagonal estimate in Section 5, the statements and proofs of Theorems 6.1–6.2, and Appendix A. The full layer-potential expansion, every estimate in Section 5, and the numerical multipole implementation were not independently reconstructed. No claim of complete proof certification or independent acceptance is made.\par
The distinction between small residual and proximity to an exact eigenvector is material. Theorem 6.1 supplies vectors with an exponentially localized profile and an explicit residual bound. Theorem 6.2 transfers a residual bound to the dense matrix under the decay hypothesis \(0<\rho<1-\varepsilon _1\). Its proof uses the action of the omitted tails on a localized vector, rather than asserting that the entire discarded matrix has small norm. The \(1/|i-j|\) entry bound alone would not justify the latter assertion. The geometric-series denominators in (6.7) are positive under the stated hypothesis, and its two tail contributions match the proof on pp. 23–25.\par
\section*{Confirmed local finding}
Problem 1 (Minor: infinite-operator spectrum in Theorem A.3).  On p. 33, Theorem A.3 includes \[\Lambda(A)+\Delta_\varepsilon
 \ \subset\ \Lambda_\varepsilon(A_{N,l}),
 \qquad \varepsilon=C\rho^N,\] as one term of its asserted union, for every \(r<1\), \(\rho>r\), with \(C\) independent of \(N\). Here \(A\) is the semi-infinite Toeplitz operator, \(A_{N,l}\) its finite section, and \(\Delta_\varepsilon\) the open disk about zero. Pages 31–32 distinguish this operator from a Laurent operator and define \(\Lambda(A)\) as its spectrum. The inclusion above is false.\par
Assumptions and decisive counterexample. On \(\ell^2(\mathbb N)\), take \[(Au)_j=u_{j+1},\qquad
 a_1=1,\quad a_k=0\ (k\ne1),\qquad l=1,\quad f(z)=z.\] This is a bounded Toeplitz operator satisfying the exact bandwidth condition \(a_k=0\) for \(|k|>l\). Set \(r=1/4\), \(\rho=1/2\), and \(\lambda=3/4\). All radius hypotheses hold, including \(\rho<1\) if that additional restriction were imposed. The nonzero square-summable vector \((1,\lambda,\lambda^2,\ldots)\) is an eigenvector of \(A\). Consequently \(\lambda\in\Lambda(A)+\Delta_\varepsilon\) for every \(\varepsilon>0\), since \(0\in\Delta_\varepsilon\).\par
Equation reference: eq:resolvent. Let \(J_N=A_{N,1}\). It has ones on the first superdiagonal and \(J_N^N=0\). Since \(\lambda\ne0\), \[(\lambda I-J_N)^{-1}
   =\sum_{k=0}^{N-1}\lambda^{-k-1}J_N^k,\qquad
 \|(\lambda I-J_N)^{-1}\|_2
   \le 4\big((4/3)^N-1\big).
 \] For every fixed finite \(C>0\), therefore, \[C2^{-N}\|(\lambda I-J_N)^{-1}\|_2
 <4C(2/3)^N\longrightarrow0.\] For sufficiently large \(N\), the resolvent norm is strictly less than \((C2^{-N})^{-1}\). Also \(\lambda\notin\Lambda(J_N)=\{0\}\). Thus \(\lambda\notin\Lambda_{C2^{-N}}(J_N)\), contradicting Theorem A.3 by Definition A.2(iii). Equivalently, every unit vector has residual strictly greater than \(C2^{-N}\). The strict inequality makes the convention at the pseudospectral boundary immaterial. Allowing \(C\) to depend on this fixed \(\lambda\) does not repair the contradiction.\par
Counterevidence and impact. The truncated geometric vector has residual of order \(\lambda^N\); the example does not dispute exponential localization at an appropriate spectral-point-dependent rate. It disputes inclusion of the entire semi-infinite spectrum at the prescribed, faster scale. It is an example within the general hypotheses of A.3, not a proposed physical gauge capacitance matrix.\par
Theorem A.4 on p. 34 omits the offending spectral-thickening term. Its proof constructs vectors from roots inside the smaller circle, cancels the boundary rows, and estimates the remaining residual. Theorem 6.1 explicitly invokes A.4 on p. 19, and Theorem 6.2 uses the entry-decay estimate from Proposition 5.9. This inspected dependency path does not use the disproved term. The Minor classification rests on that limited impact, not on the ease of a possible typographical repair.\par
Repair boundary. Replacing this particular term by \(\Lambda(A_{N,l})+\Delta_\varepsilon\) makes that term’s inclusion valid for every finite matrix: for a unit eigenvector \(w\) with eigenvalue \(\mu\), \(\|(A_{N,l}-\lambda I)w\|=|\mu-\lambda|<\varepsilon\). This verifies the replacement term only. The rest of the theorem and all uniformity claims require their own proof. The cited Reichel–Trefethen paper, Theorem 3.2, printed p. 173, was also inspected and displays the same notation; it does not remove the contradiction above. No correction adopted by the authors has been established.\par
\section*{Interpretation, computation and limitations}
The discussion following Theorem 6.2 (p. 25), the abstract, and the conclusion use language about exact eigenmodes more broadly than the residual estimates alone establish. Section 6 also acknowledges the difficulty of proving exact-eigenvector decay for dense Toeplitz matrices and uses numerical evidence. The present review does not turn this difference in scope into a second confirmed defect: it supplies no counterexample within the physical family and no proof that the claimed physical localization is false. The rigorous result verified here should be read as a pseudo-eigenvector result unless additional eigenvector stability estimates are supplied.\par
The complete 63-file repository snapshot `jinghaocao/skin\_effect`, commit `1bc8d74162f8d647b4448ea5190cd5ef71168037`, was inventoried. Static reading covered `Run\_SKIN.m`, `symbol.m`, and `produce\_condensation.m`. The code uses an inclusive diagonal cutoff whereas (6.1) uses \(|i-j|<k\); this is a parameter-convention distinction, with no demonstrated numerical error here. The scripts use real parts of computed eigenvectors in some plots; this review did not establish that this changes any displayed conclusion. Author code was not executed, and computed figures were not reproduced.\par
The recorded scientific text coverage is pp. 1–5, 9, 13–14, 16–25, 31–37 and 39; selected evidence pages were also viewed as images. Uninspected technical passages and numerical data remain limitations. Source reading, impact judgment, report composition, compilation, and page inspection were performed in the present review conversation. The administrative scientific decision remains pending, with no human review or independent acceptance recorded.\par
9 H. Ammari, S. Barandun, J. Cao, B. Davies, E. O. Hiltunen and P. Liu, The non-Hermitian skin effect with three-dimensional long-range coupling, J. Eur. Math. Soc., online first, 4 August 2025, doi:10.4171/JEMS/1685 (https://doi.org/10.4171/JEMS/1685). Exact supplied 39-page CaltechAUTHORS copy. L. Reichel and L. N. Trefethen, Eigenvalues and pseudo-eigenvalues of Toeplitz matrices, Linear Algebra Appl. 162–164 (1992), 153–185, doi:10.1016/0024-3795(92)90374-J (https://doi.org/10.1016/0024-3795(92)90374-J). Theorem 3.2 consulted in the author’s Oxford-hosted copy.\par
\end{document}
